\chapter{Categorical Proofs of LSP Modalities}
\label{apx:modalities}

This appendix provides the consolidated mathematical proofs required to establish Homological Boundedness across the distinct operational classes of Language Server Protocol (LSP) execution end-points under stochastic load. Rather than enumerating isomorphic proofs for individual endpoints, we partition the protocol surface into three fundamental topological classes: Read-Only Modalities, Mutational Modalities, and Global Queries.

\section{Class I: Read-Only Modalities (Presheaf Resolution)}

This class encompasses operations that query the state manifold $\mathcal{M}$ without injecting perturbation vectors. Examples include \texttt{textDocument/hover}, \texttt{textDocument/semanticTokens}, and \texttt{textDocument/signatureHelp}.

Let the stochastic request vector for a Class I operation be denoted as $\mathbf{X}_{read}$. Because these operations do not mutate the AST, the covariant derivative along the Levi-Civita connection evaluates strictly to zero:
\begin{equation}
\nabla_{X_{read}} Y = 0 \quad \forall Y \in T\mathcal{M}
\end{equation}
The computational cost is therefore entirely restricted to evaluating the Yoneda embedding $H_A = \text{Hom}_{\mathbf{AST}}(-, A)$. For a hover request at coordinate $q_0$, the engine accesses the localized presheaf data $F(q_0)$ within the \texttt{salsa} incremental database. 

Because the evaluation requires zero modification to the Christoffel symbols ($\Delta \Gamma^k_{ij} = 0$), the operational work is bounded by the hash-map lookup latency:
\begin{equation}
\int_{q(\tau)}^{q(\tau)} g_{ij} \frac{dq^i}{d\tau} \frac{dq^j}{d\tau} d\tau = 0
\end{equation}
This mathematically guarantees that the swarm cannot induce semantic backpressure through Read-Only bombardment, provided the $C^*$-algebraic boundary filters invalid JSON-RPC schemas in $\mathcal{O}(1)$ prior to query resolution.

\section{Class II: Mutational Modalities (Differential Perturbation)}

This class encompasses state-altering operations, primarily \texttt{textDocument/didChange} and \texttt{textDocument/codeAction}.

Let the mutational vector be $\mathbf{X}_{mut}$. As proven in Chapter \ref{ch:differential_calculus}, if $\mathbf{X}_{mut} \in E_B^*$, the request permeates to the \texttt{salsa} graph. The engine must compute the differential perturbation over the Riemannian manifold $\mathcal{M}$. 

By calculating the specific Christoffel symbols $\Gamma^{k}_{ij}$ localized solely to the topological domain affected by $\mathbf{X}_{mut}$, the integration of the covariant derivative $\nabla$ over the tangent space guarantees that:
\begin{equation}
\int_{q(\tau)}^{q(\tau+\Delta)} g_{ij} \frac{dq^i}{d\tau} \frac{dq^j}{d\tau} d\tau \ll \oint_{\mathcal{M}} \mathcal{O}(N) d\mathcal{M}
\end{equation}
This bounds the latency of the AST transformation strictly to the neighborhood of the edit $|\Delta|$, successfully isolating the system from the global document size $N$.

\section{Class III: Global Queries (Geodesic Traversal)}

This class encompasses operations that necessitate exploring the broader interconnected graph of the repository, such as \texttt{workspace/symbol} or \texttt{textDocument/references}.

Let the global query vector be $\mathbf{X}_{global}$. These requests represent the highest computational risk, as naive implementations degrade to full tree traversals. In the \texttt{tower-lsp-max} architecture, these queries are resolved not via tree traversal, but via contraction of the metric tensor over pre-computed indexed sub-manifolds.

Let $R_{ij}$ be the Ricci curvature tensor representing the density of inter-node dependencies. The \texttt{AutoLspAdapter} maintains a continuous projection of this curvature within the \texttt{salsa} cache. The work to resolve $\mathbf{X}_{global}$ is bounded by tracing the geodesics generated by the pre-computed connections:
\begin{equation}
W_{global} \propto \text{Tr}(R_{ij} X_{global}^i X_{global}^j)
\end{equation}
Because the Ricci curvature is maintained incrementally during Class II operations, the query execution for Class III avoids divergent infinities, resolving in time proportional to the frequency of the requested symbol rather than the volume of the repository.